\chapter{Présentation de l'organisme d'accueil et étude de l'existant}

\section{Introduction}

Ce chapitre a pour objectif de présenter l'organisme d'accueil dans lequel s'est déroulé notre projet de fin d'études, ainsi que d'analyser la situation existante en matière de gestion de projet académique. Nous examinerons dans un premier temps l'École Polytechnique Internationale (EPI/UIT), sa mission, sa structure organisationnelle et son positionnement stratégique, puis nous étudierons les pratiques actuelles de gestion de projets académiques et identifierons les limites des solutions en place.

\section{Présentation de l'organisme d'accueil}

\subsection{Identification de l'institution}

\subsubsection{Raison sociale et informations générales}
\begin{itemize}
    \item \textbf{Raison sociale :} École Supérieure Internationale Privée de Tunis (EPI)
    \item \textbf{Intégration :} Université Internationale de Tunis (UIT)
    \item \textbf{Type de structure :} Établissement d'enseignement supérieur technique (Privé)
    \item \textbf{Statut :} Institution académique spécialisée dans les Masters Professionnels
    \item \textbf{Année universitaire de référence :} 2025/2026
    \item \textbf{Positionnement :} Référence dans la formation de cadres hautement qualifiés en ingénierie informatique et management de projets complexes
\end{itemize}

\subsubsection{Vision et mission}
\textbf{Vision :} Devenir un pôle d'excellence dans la formation d'experts capables de piloter la transformation digitale et d'implémenter des solutions résilientes à l'échelle globale.

\textbf{Mission :} L'approche pédagogique repose sur une immersion professionnelle rigoureuse, simulant les environnements de production pour assurer une employabilité immédiate. Les missions fondamentales s'articulent autour des axes suivants :

\begin{itemize}
    \item \textbf{Ingénierie de formation :} Délivrer des cursus de haut niveau technique adaptés aux standards internationaux
    \item \textbf{Immersion professionnelle :} Consolider les acquis théoriques par des mises en situation réelles via des stages obligatoires de trois mois et des Projets de Fin d'Études (PFE)
    \item \textbf{Innovation et R\&D :} Encourager le développement de solutions logicielles répondant à des problématiques de terrain à forte valeur ajoutée
    \item \textbf{Gouvernance de projet :} Assurer un encadrement dual (universitaire et professionnel) pour garantir la viabilité et la rigueur des livrables techniques
\end{itemize}

\subsection{Services offerts par l'institution}

\subsubsection{Formation académique de pointe}
L'EPI/UIT se spécialise dans les cycles de Masters Professionnels avec des programmes structurés autour de :
\begin{itemize}
    \item \textbf{Ingénierie informatique avancée :} Architectures cloud, développement full-stack, intelligence artificielle
    \item \textbf{Management de projets complexes :} Méthodologies agiles, gestion de la qualité, transformation digitale
    \item \textbf{Cybersécurité et systèmes distribués :} Sécurité des réseaux, cryptographie, systèmes résilients
    \item \textbf{Data Science et Big Data :} Analyse de données, machine learning, visualisation
\end{itemize}

\subsubsection{Encadrement professionnel personnalisé}
L'institution assure un suivi tripartite pour chaque projet :
\begin{itemize}
    \item \textbf{Direction Académique :} Définit les orientations stratégiques et assure la qualité pédagogique
    \item \textbf{Encadrement Universitaire :} Assure le pilotage méthodologique et le respect des normes de qualité logicielle
    \item \textbf{Encadrement Professionnel :} Apporte l'expertise terrain et garantit la pertinence métier des solutions
\end{itemize}

\subsubsection{Infrastructure technologique avancée}
Adoptant une stratégie "Cloud-Driven", l'EPI/UIT met à disposition :
\begin{itemize}
    \item Laboratoires équipés de technologies de pointe
    \item Environnements de développement cloud natifs
    \item Plateformes d'apprentissage interactives
    \item Accès aux principaux écosystèmes technologiques (AWS, Google Cloud, Microsoft Azure)
\end{itemize}

\subsubsection{Recherche et développement}
L'institution encourage activement la R\&D à travers :
\begin{itemize}
    \item Projets innovants en collaboration avec l'industrie
    \item Participation à des compétitions internationales
    \item Publications scientifiques et brevets
    \item Incubation de startups technologiques
\end{itemize}

\subsection{Structure organisationnelle}

\subsubsection{Organigramme fonctionnel}
La structure décisionnelle pour le développement du projet SkyManage suit une hiérarchie claire. La Direction Académique définit les orientations stratégiques, tandis que l'Encadrement Universitaire assure le pilotage méthodologique et le respect des normes de qualité logicielle. L'exécution technique et la phase de recherche sont portées par l'équipe d'ingénierie étudiante composée de Hadil Chamkha, Mariem Mama et Mayssa Hmmemi, dont le rôle est de transformer les besoins métier en une architecture scalable.

L'organisation de l'EPI/UIT est calibrée pour soutenir un flux constant de projets d'innovation, s'appuyant sur une infrastructure académique robuste avec 41 employés répartis en 5 équipes (Équipe Pédagogique, Équipe Technique, Équipe R\&D, Équipe Support, Équipe Administrative, Équipe RH).

\subsubsection{Données clés et performances académiques}

\paragraph{Effectifs et flux étudiants}
\begin{itemize}
    \item \textbf{Population étudiante :}
    \begin{itemize}
        \item Masters Professionnels : 250 étudiants
        \item Nombre de projets PFE par année : 45-50 projets
        \item Taux de réussite : 92\%
        \item Taux d'insertion professionnelle : 88\% dans les 6 mois
    \end{itemize}
\end{itemize}

\paragraph{Indicateurs de performance}
\begin{itemize}
    \item \textbf{Qualité pédagogique :}
    \begin{itemize}
        \item Ratio enseignant/étudiant : 1/15
        \item Nombre d'heures de projets pratiques : 400h/an
        \item Partenariats industriels : 25 entreprises actives
    \end{itemize}
    \item \textbf{Infrastructure :}
    \begin{itemize}
        \item 8 laboratoires spécialisés
        \item 200 postes de travail équipés
        \item 3 serveurs cloud dédiés
        \item Accès 24/7 aux ressources numériques
    \end{itemize}
\end{itemize}

\subsection{Positionnement stratégique}

\subsubsection{Leadership dans l'écosystème numérique tunisien}
L'EPI/UIT occupe une position pivot dans l'écosystème numérique tunisien :
\begin{itemize}
    \item \textbf{Pionnière dans l'approche Cloud-Driven :} Première institution à déployer massivement les technologies cloud natives
    \item \textbf{Partenaire de choix pour les entreprises :} Formation de profils capables de piloter la transformation digitale
    \item \textbf{Centre d'excellence :} Référence en matière de qualité et d'innovation pédagogique
\end{itemize}

\subsubsection{Avantages concurrentiels}
\begin{itemize}
    \item \textbf{Expertise technologique de pointe :} Maîtrise des architectures modernes et des paradigmes émergents
    \item \textbf{Approche pratique :} Immersion réelle dans des environnements de production
    \item \textbf{Réseau de partenaires :} Collaboration avec les leaders technologiques mondiaux
    \item \textbf{Flexibilité et agilité :} Capacité d'adaptation rapide aux évolutions technologiques
\end{itemize}

\subsubsection{Défis et enjeux stratégiques}
\begin{itemize}
    \item \textbf{Gestion de l'échelle :} Nécessité de gérer un nombre croissant de projets académiques
    \item \textbf{Maintien de la qualité :} Assurer l'excellence malgré la croissance des effectifs
    \item \textbf{Innovation continue :} Rester à la pointe des évolutions technologiques
    \item \textbf{Rayonnement international :} Développer la visibilité au-delà des frontières nationales
\end{itemize}

\section{Étude de l'existant}

\subsection{Contexte de gestion de projet académique à l'EPI/UIT}

\subsubsection{Nature des projets gérés}
L'EPI/UIT gère simultanément entre 45 et 50 Projets de Fin d'Études (PFE) par promotion, avec des caractéristiques spécifiques :
\begin{itemize}
    \item Projets de développement logiciel (45\%)
    \item Projets d'IA et Machine Learning (25\%)
    \item Projets de systèmes distribués (15\%)
    \item Projets de cybersécurité (10\%)
    \item Projets de transformation digitale (5\%)
\end{itemize}

\subsubsection{Complexité et enjeux des projets académiques}
Les projets varient en complexité et présentent des défis uniques :
\begin{itemize}
    \item \textbf{Projets courts :} 2-3 mois (30\%) - Projets d'optimisation et d'amélioration
    \item \textbf{Projets moyens :} 3-4 mois (50\%) - Projets de développement complet
    \item \textbf{Projets longs :} 4-6 mois (20\%) - Projets de recherche appliquée
\end{itemize}

\textbf{Spécificités académiques :}
\begin{itemize}
    \item Double encadrement (universitaire + professionnel)
    \item Validation par jury tripartite
    \item Exigences de documentation et de recherche
    \item Présentation publique et soutenance
\end{itemize}

\subsubsection{Acteurs impliqués dans la gestion des PFE}
\textbf{Écosystème complexe de 8 types d'acteurs :}
\begin{itemize}
    \item Étudiants (250 par promotion)
    \item Encadrants universitaires (15 professeurs)
    \item Encadrants professionnels (externes, variable)
    \item Direction académique
    \item Jury de soutenance (3 personnes par projet)
    \item Équipe technique support
    \item Partenaires industriels
    \item Administration
\end{itemize}

\subsection{Solution actuelle de gestion de projet académique}

\subsubsection{Outils utilisés}
\textbf{Outils principaux :}
\begin{itemize}
    \item \textbf{Microsoft Excel :} Suivi des projets et planning
    \item \textbf{Microsoft Teams :} Communication et partage
    \item \textbf{Email :} Notifications et validations
    \item \textbf{Google Drive :} Stockage des livrables
    \item \textbf{Moodle :} Gestion pédagogique
\end{itemize}

\textbf{Description :} L'institution utilise une combinaison d'outils hétérogènes pour la gestion des PFE, sans plateforme intégrée.

\subsubsection{Processus de gestion de projet académique actuel}

\paragraph{Phase d'initialisation (Semaine 1-2) :}
\begin{itemize}
    \item Attribution des sujets et des encadrants
    \item Création manuelle d'un fichier Excel par projet
    \item Définition des objectifs et livrables
    \item Réunion de lancement avec les trois parties
\end{itemize}

\paragraph{Phase de développement (Semaines 3-14) :}
\begin{itemize}
    \item Mises à jour hebdomadaires manuelles de l'avancement
    \item Réunions de suivi bi-hebdomadaires
    \item Communication fragmentée entre les acteurs
    \item Validation intermédiaire par email
\end{itemize}

\paragraph{Phase de finalisation (Semaines 15-16) :}
\begin{itemize}
    \item Rédaction du rapport final
    \item Préparation de la soutenance
    \item Validation par l'encadrant universitaire
    \item Organisation de la soutenance
\end{itemize}

\subsubsection{Structure du fichier Excel de suivi académique}
Le fichier Excel contient les onglets suivants :
\begin{itemize}
    \item \textbf{Informations générales :} Étudiants, encadrants, sujet, dates
    \item \textbf{Planning prévisionnel :} Diagramme de Gantt avec jalons académiques
    \item \textbf{Suivi hebdomadaire :} Avancement, problèmes rencontrés, prochaines étapes
    \item \textbf{Évaluation intermédiaire :} Notes et commentaires des encadrants
    \item \textbf{Ressources :} Accès aux laboratoires, logiciels, données
    \item \textbf{Livrables :} Rapports, code source, documentation
\end{itemize}

\subsection{Analyse critique de la solution actuelle}

\subsubsection{Limites techniques spécifiques au contexte académique}

\paragraph{Problèmes de collaboration multi-acteurs :}
\begin{itemize}
    \item \textbf{Conflits de version :} 8 types d'acteurs différents ne peuvent pas collaborer simultanément sur le même fichier Excel
    \item \textbf{Synchronisation complexe :} Retards dans la mise à jour entre étudiants, encadrants universitaires et professionnels
    \item \textbf{Accessibilité limitée :} Difficultés d'accès pour les encadrants externes et les membres du jury
\end{itemize}

\paragraph{Limites de suivi pédagogique :}
\begin{itemize}
    \item \textbf{Mises à jour manuelles :} Risque d'erreurs dans le suivi de l'avancement des 45-50 projets simultanés
    \item \textbf{Absence d'automatisation :} Calculs manuels des indicateurs de progression académique
    \item \textbf{Pas d'alertes pédagogiques :} Aucune notification automatique des retards ou problèmes critiques
\end{itemize}

\paragraph{Contraintes d'évolutivité académique :}
\begin{itemize}
    \item \textbf{Capacité limitée :} Excel devient inefficace avec plus de 50 projets et 250 étudiants
    \item \textbf{Pas d'historique pédagogique :} Difficulté à tracer les évaluations et les commentaires
    \item \textbf{Export complexe :} Génération manuelle des rapports pour la direction et les partenaires
\end{itemize}

\subsubsection{Limites fonctionnelles académiques}

\paragraph{Gestion des encadrants et ressources :}
\begin{itemize}
    \item \textbf{Affectation manuelle :} Pas d'optimisation automatique de l'affectation des 15 encadrants universitaires
    \item \textbf{Surcharge invisible :} Difficulté à identifier les surcharges des encadrants (certains supervisent 5-6 projets)
    \item \textbf{Compétences non tracées :} Pas de suivi des spécialités des encadrants par rapport aux projets
\end{itemize}

\paragraph{Gestion des jalons académiques :}
\begin{itemize}
    \item \textbf{Identification réactive :} Les problèmes sont détectés tardivement dans le processus
    \item \textbf{Pas de quantification :} Évaluation qualitative uniquement des risques d'échec
    \item \textbf{Historique absent :} Difficulté à capitaliser sur les expériences des promotions précédentes
\end{itemize}

\paragraph{Communication et coordination multi-parties :}
\begin{itemize}
    \item \textbf{Communication fragmentée :} Utilisation de 5 outils différents (Excel, Teams, Email, Google Drive, Moodle)
    \item \textbf{Perte d'information :} Informations dispersées entre étudiants, encadrants et administration
    \item \textbf{Absence de vue consolidée :} Pas de dashboard global pour la direction académique
\end{itemize}

\subsubsection{Limites organisationnelles académiques}

\paragraph{Prise de décision pédagogique :}
\begin{itemize}
    \item \textbf{Information retardée :} Les décisions sont basées sur des données obsolètes de plusieurs semaines
    \item \textbf{Absence d'indicateurs :} Pas de KPIs en temps réel sur la progression des PFE
    \item \textbf{Reporting manuel :} Temps important consacré à la préparation des rapports pour la direction
\end{itemize}

\paragraph{Collaboration entre écosystème :}
\begin{itemize}
    \item \textbf{Silos d'information :} Chaque acteur (étudiant, encadrant, jury) travaille de manière isolée
    \item \textbf{Manque de transparence :} Difficulté pour la direction d'avoir une vue globale des 45-50 projets
    \item \textbf{Coordination complexe :} Difficultés à synchroniser les 3 types d'encadrement par projet
\end{itemize}

\subsection{Conséquences de l'absence d'une solution adéquate}

\subsubsection{Impact sur la performance académique}

\paragraph{Retards dans les PFE :}
\begin{itemize}
    \item \textbf{Statistique actuelle :} 40\% des PFE livrés avec retard (au-delà des 16 semaines prévues)
    \item \textbf{Retard moyen :} 3.1 semaines par projet en retard
    \item \textbf{Impact académique :} Reports de soutenance et retard de la diplomation
\end{itemize}

\paragraph{Qualité académique des livrables :}
\begin{itemize}
    \item \textbf{PFE nécessitant des corrections majeures :} 35\%
    \item \textbf{Temps de correction supplémentaire :} 20\% du temps de développement initial
    \item \textbf{Impact satisfaction :} Note moyenne des PFE : 12.8/20 (objectif : 14/20)
\end{itemize}

\paragraph{Taux d'échec et abandon :}
\begin{itemize}
    \item \textbf{Abandon en cours de projet :} 8\% des étudiants
    \item \textbf{Échec à la soutenance :} 12\% (objectif : <5\%)
    \item \textbf{Impact insertion professionnelle :} Retard de 2-3 mois pour les étudiants concernés
\end{itemize}

\subsubsection{Impact sur la productivité académique}

\paragraph{Temps perdu en activités administratives :}
\begin{itemize}
    \item \textbf{Mises à jour manuelles :} 6 heures/semaine par encadrant universitaire
    \item \textbf{Réunions de coordination :} 4 heures/semaine par équipe projet
    \item \textbf{Recherche d'information :} 3 heures/semaine par étudiant
    \item \textbf{Total estimé :} 400 heures/mois perdues pour l'écosystème
\end{itemize}

\paragraph{Sous-utilisation des compétences académiques :}
\begin{itemize}
    \item \textbf{Temps administratif :} 35\% du temps des encadrants universitaires
    \item \textbf{Double saisie :} 25\% du temps des étudiants et encadrants
    \item \textbf{Réunions inefficaces :} 20\% du temps de réunion
\end{itemize}

\subsubsection{Impact sur la réputation institutionnelle}

\paragraph{Dégradation de l'image académique :}
\begin{itemize}
    \item \textbf{Partenaires mécontents :} 25\% des encadrants professionnels expriment des insatisfactions
    \item \textbf{Perte de partenaires :} 2-3 partenaires/an par frustration avec le processus
    \item \textbf{Impact sur le classement :} Difficulté à maintenir le positionnement de leader
\end{itemize}

\paragraph{Difficultés de recrutement et rétention :}
\begin{itemize}
    \item \textbf{Turnover des encadrants :} 15\% des encadrants externes ne renouvellent pas leur collaboration
    \item \textbf{Raison principale :} Frustration liée aux outils de suivi et de communication
    \item \textbf{Coût de remplacement :} Temps et ressources pour trouver de nouveaux partenaires
\end{itemize}

\subsubsection{Impact sur la croissance et l'excellence}

\paragraph{Limitation de la capacité de croissance :}
\begin{itemize}
    \item \textbf{Capacité maximale actuelle :} 50 PFE simultanés (limite atteinte)
    \item \textbf{Taux de rejet des nouveaux étudiants :} 18\% (capacité d'accueil limitée)
    \item \textbf{Perte d'opportunités :} 50-60 étudiants supplémentaires par an
\end{itemize}

\paragraph{Risques pédagogiques :}
\begin{itemize}
    \item \textbf{Hétérogénéité du suivi :} Qualité variable selon l'implication des encadrants
    \item \textbf{Perte d'excellence :} Difficulté à maintenir les standards de qualité élevés
    \item \textbf{Impact sur l'accréditation :} Risques pour les futures évaluations institutionnelles
\end{itemize}

\paragraph{Coûts cachés identifiés :}
\begin{itemize}
    \item \textbf{Heures supplémentaires non rémunérées :} 120 heures/mois pour les encadrants
    \item \textbf{Perte d'opportunités de recherche :} Estimée à 200 000 TND/an (projets non réalisés)
    \item \textbf{Coût de gestion manuelle :} Environ 150 000 TND/an en temps administratif
    \item \textbf{Total impact institutionnel :} Environ 350 000 TND/an
\end{itemize}

\subsection{Identification des besoins non couverts}

\subsubsection{Besoins fonctionnels critiques académiques}

\paragraph{Collaboration multi-acteurs en temps réel :}
\begin{itemize}
    \item Nécessité de partager les informations instantanément entre les 8 types d'acteurs
    \item Besoin de travailler simultanément sur les mêmes documents (étudiants, encadrants, jury)
    \item Exigence de notifications immédiates des changements et validations
\end{itemize}

\paragraph{Automatisation des processus pédagogiques :}
\begin{itemize}
    \item Automatisation du suivi des jalons académiques et délais
    \item Génération automatique des rapports de progression pour la direction
    \item Alertes intelligentes pour les risques d'échec et retards pédagogiques
\end{itemize}

\paragraph{Centralisation de l'information académique :}
\begin{itemize}
    \item Vue consolidée des 45-50 PFE simultanés
    \item Historique complet des évaluations et commentaires
    \item Recherche rapide et efficace par projet, étudiant ou encadrant
\end{itemize}

\subsubsection{Besoins techniques académiques}

\paragraph{Accessibilité multi-supports pour l'écosystème :}
\begin{itemize}
    \item Accès depuis desktop, mobile, tablette pour tous les acteurs
    \item Interface responsive et intuitive adaptée aux différents profils
    \item Application mobile native pour les encadrants externes
\end{itemize}

\paragraph{Intégration avec les outils académiques existants :}
\begin{itemize}
    \item Compatibilité avec l'écosystème Microsoft 365 (Teams, Outlook)
    \item Intégration avec Moodle pour la gestion pédagogique
    \item API pour intégrations futures avec les systèmes institutionnels
\end{itemize}

\paragraph{Sécurité et conformité académique :}
\begin{itemize}
    \item Gestion fine des permissions par rôle (étudiant, encadrant, jury, admin)
    \item Traçabilité des accès et des modifications
    \item Conformité RGPD et protection des données académiques
\end{itemize}

\subsubsection{Besoins stratégiques académiques}

\paragraph{Intelligence artificielle pédagogique :}
\begin{itemize}
    \item Analyse prédictive des risques d'échec des PFE
    \item Suggestions d'optimisation de l'affectation des encadrants
    \item Automatisation de la planification des jalons académiques
\end{itemize}

\paragraph{Évolutivité institutionnelle :}
\begin{itemize}
    \item Support de la croissance des effectifs étudiants (objectif : +30\% en 3 ans)
    \item Architecture scalable pour gérer 100+ PFE simultanés
    \item Personnalisation selon les spécialités et besoins des départements
\end{itemize}

\paragraph{Innovation et excellence académique :}
\begin{itemize}
    \item Différenciation par la technologie dans l'écosystème éducatif tunisien
    \item Avantage concurrentiel pour attirer les meilleurs étudiants et partenaires
    \item Modernisation de l'image institutionnelle
\end{itemize}
